\documentclass[12pt]{article}

\usepackage[a4paper, margin=1in]{geometry}
\usepackage{amsmath, amssymb}
\usepackage{setspace}
\usepackage{hyperref}

\setstretch{1.2}

\title{\textbf{AgenticBrain: \\ The Global Orchestrated Decision (GOD-BRAIN) Framework for Autonomous Neural Agency}}
\author{Iyer Ramkumar Ramasubramanian \\ \textit{Independent Researcher}}
\date{\today}

\begin{document}

\maketitle

\begin{abstract}
While previous models like the Homo Deus Brain Model (HDBM) demonstrate system stability and passive equilibrium, they often prioritize stability over utility in the absence of external objectives. This work introduces the \textbf{AgenticBrain}, utilizing the \textbf{Global Orchestrated Decision (GOD-BRAIN)} framework. This architecture enables an agentic neural model to move beyond reactive processing into proactive environmental orchestration. By integrating training logs from agentic simulations, we demonstrate a system capable of balancing biological constraints with complex, multi-stage objective functions.
\end{abstract}

\section{Introduction}
Current advancements in neural architectures are shifting from static inference to agentic autonomy. Traditional AI systems operate within isolated feedback loops that often lead to "passive equilibrium"—a state where minimal action is preferred to conserve energy. The AgenticBrain addresses this limitation by implementing a high-level orchestration layer that allows the neural system to act as an independent professional entity within its environment.

\section{The GOD-BRAIN Framework}

The Global Orchestrated Decision (GOD-BRAIN) framework is built upon four architectural pillars:

\begin{itemize}
\item \textbf{Proactive Agency}: Unlike reactive models, the AgenticBrain generates internal objectives that drive action even when environmental stimuli are neutral.
\item \textbf{Global Orchestration}: Decisions are not localized to specific sensors but are orchestrated across the entire neural substrate, including the IoT and Nano layers.
\item \textbf{Dynamic Utility}: The system balances the biological cost (energy and stress) against the utility of achieving high-level goals.
\item \textbf{Scale-Invariant Decisioning}: The framework is designed to function from individual neural simulations to national-scale infrastructures.
\end{itemize}

\section{System Architecture}

The AgenticBrain preserves the closed-loop feedback of the HDBM but adds an \textbf{Orchestration Layer}:

\subsection{Orchestration Layer}
Acts as the "Global Brain," synthesizing data from the Cognitive and Biological layers to prioritize long-term goals over immediate energy conservation.

\subsection{Agentic Actuator Layer}
Executes complex, multi-stage sequences in the environment, such as managing a GitHub repository or generating research documentation.

\section{Mathematical Formulation}

In the AgenticBrain, the state transition includes a "Goal Vector" $G$:

\[
X_{t+1} = f(WX_t + S(X_t) + B(X_t) + G_t)
\]

Where $G_t$ represents the orchestrated objective. The loss function is modified to include a "Utility Gain" $U$:

\[
L = ||A_t||^2 - \alpha(U_t)
\]

By adjusting the coefficient $\alpha$, the system can be tuned to prioritize goal achievement even when the physical action magnitude $||A_t||^2$ is high, preventing the system from falling into the "passive equilibrium" trap observed in earlier models.

\section{Results and Simulation Logs}

Experimental training logs for the AgenticBrain indicate:
\begin{itemize}
\item \textbf{Goal Persistence}: The system maintains action sequences despite increasing biological stress variables.
\item \textbf{Optimized Equilibrium}: Instead of zero-action, the system reaches a "Active Equilibrium" where energy consumption is high but utility output is maximized.
\item \textbf{Self-Correction}: The agentic layer identifies and overrides sub-optimal "passive" weights to ensure task completion.
\end{itemize}

\section{Conclusion}
The AgenticBrain and the GOD-BRAIN framework provide a blueprint for post-human intelligence systems that possess true agency. By treating the neural network as a global orchestrator, we move closer to a unified theory of neurobiology and artificial decision-making.

\section{Repository for Experimentation}
\noindent
\url{https://github.com/spacetracker-collab/AgenticBrain}

\end{document}
